\newcommand{\boolNo}{\text{No}} % boolean No
\newcommand{\boolYes}{\text{Yes}} % boolean Yes

\NewDocumentCommand{\localG}{g}{%
    \IfNoValueTF{#1}
    {G}%
    {G_{#1}}%
}
\NewDocumentCommand{\sparseG}{g}{%
    \IfNoValueTF{#1}
    {\widetilde{G}}%
    {\widetilde{G}_{#1}}%
} % sparsified graph in dynamic agreement (DA)
\NewDocumentCommand{\dyCC}{g}{%
    \IfNoValueTF{#1}
    {C}%
    {C_{#1}}%
} % dynamic data structure for connected components
\NewDocumentCommand{\journal}{g}{%
    \IfNoValueTF{#1}
    {\mathcal{J}}%
    {\mathcal{J}_{#1}}%
}
\newcommand{\epsDA}{\varepsilon_{\rm DA}} % accuracy parameter for dynamic agreement
\NewDocumentCommand{\interestSet}{g}{%
    \IfNoValueTF{#1}
    {\mathcal{I}}%
    {\mathcal{I}_{#1}}%
} % interesting set in DA
\NewDocumentCommand{\anchorSet}{g}{%
    \IfNoValueTF{#1}
    {\Phi}%
    {\Phi_{#1}}%
} % anchor set in DA


\newcommand{\openN}[2]{N_{#1}(#2)} % open neighborhood of vertex #2 in graph #1
\newcommand{\closedN}[2]{N_{#1}[#2]} % open neighborhood of vertex #2 in graph #1
\newcommand{\snapshotN}[1]{\widetilde{N}(#1)} % open neighborhood of vertex #2 in graph #1


\newcommand{\explicitDA}{\mathrm{ExplicitDynamicAgreement}} % explicit dynamic agreement algorithm
\NewDocumentCommand{\Notify}{g}{%
    \IfNoValueTF{#1}
    {\mathrm{Notify}}%
    {\mathrm{Notify}(#1)}%
}
\NewDocumentCommand{\Anchor}{g}{%
    \IfNoValueTF{#1}
    {\mathrm{Anchor}}%
    {\mathrm{Anchor}(#1)}%
}
\NewDocumentCommand{\Clean}{g}{%
    \IfNoValueTF{#1}
    {\mathrm{Clean}}%
    {\mathrm{Clean}(#1)}%
}
\NewDocumentCommand{\Connect}{g}{%
    \IfNoValueTF{#1}
    {\mathrm{Connect}}%
    {\mathrm{Connect}(#1)}%
}
\NewDocumentCommand{\LabelSync}{g}{%
    \IfNoValueTF{#1}
    {\mathrm{LabelSync}}%
    {\mathrm{LabelSync}(#1)}%
}

\newcommand{\delt}[1]{\mathrm{Del}_{#1}}
\newcommand{\addt}[1]{\mathrm{Add}_{#1}}
\newcommand{\delIncEdge}[1]{\delta(#1)} % number of edges incident on vertex #1 that are deleted (\in \delt)
\NewDocumentCommand{\ccclass}{g}{%
    \IfNoValueTF{#1}
    {C}%
    {C_{#1}}%
} % connected component class
\NewDocumentCommand{\ccclasses}{g}{%
    \IfNoValueTF{#1}
    {\mathcal C}%
    {\mathcal{C}_{#1}}%
} % connected component class
\newcommand{\comp}[1]{\mathrm{comp}(#1)} % connected component of vertex #1
\NewDocumentCommand{\piece}{ g }{%
    \IfNoValueTF{#1}
    {piece}%
    {piece#1}%
}
\newcommand{\retainedPieces}[1]{R_{#1}} % retained pieces of class #1
\newcommand{\piecesToMerge}{\mathcal{Z}} %
\newcommand{\pieceToMerge}{P} %

\newcommand{\successDA}{\mathcal{E}_{\rm DA}} % success event of dynamic agreement

\newcommand{\edSolver}{\mathcal D}
\newcommand{\edFactor}{\eta} % approximation factor of the explicit correlation-clustering algorithm
\newcommand{\edTime}[1]{T_{\mathcal D}(#1)}
\newcommand{\edcost}[1]{\widehat{\cost}_{#1}}
\newcommand{\edmean}[2]{\widehat{\cost}^{(#1)}_{#2}}
\newcommand{\edscore}[2]{\widehat f_{#1}(#2)}
\newcommand{\edhistory}[2]{\widehat G_{#1}(#2)}
\newcommand{\edminimum}[1]{\widehat M_{#1}}
\newcommand{\edproxy}[2]{Z_{#1}(#2)}
\newcommand{\edlog}[2]{\mathcal L_{#1}(#2)}
\newcommand{\edbucket}[2]{B_{#1}(#2)}
\newcommand{\edatom}[1]{B_{123}(#1)}


\section{Explicit Dynamic Algorithm for Insertion-Only Streams}\label{sec:explicit_dynamic_algorithm_for_insertion_only_stream}

We prove the following theorem.
\ExplicitDynamic*

We use the vertex-arrival model of \cref{subsec:online-model}, with the input-label access of \cref{subsec:dynamic-model}.
In particular, the input partitions and arrival order are oblivious, the horizon $ n $ is known, and any selected input label of an arriving vertex can be read in constant time.

We obtain \cref{thm:explicit.dynamic} from the following two lemmas.

\begin{lemma}\label{lem:explicit.consensus.clustering}
Suppose that $ \edSolver $ is a randomized explicit dynamic correlation-clustering algorithm in the database model~(\cref{def:query.model.label}) that, on every oblivious stream of at most $ n $ vertex insertions, has expected amortized update time $ \edTime{n}=\polylog(n+2) $ including queries and reports, and with probability at least $ 1/2 $ maintains an $ \edFactor $-approximation simultaneously at all prefixes.
Then, for every $ 0<\eps<2/(3\edFactor+5) $ and $ \zeta\in(0,1) $, there is an explicit dynamic consensus-clustering algorithm for insertion-only streams that achieves approximation ratio $ 2-2/(3\edFactor+5)+\eps $ at all prefixes simultaneously with probability at least $ 1-\zeta $.
Its expected amortized update time is
\begin{align}
    O\!\left(
        \eps^{-3}\log^2\frac{2n}{\zeta}
        \log\frac{2n}{\eps\zeta}\cdot\polylog(n+2)
    \right).
    \label{eq:explicit.consensus.time}
\end{align}
\end{lemma}

\begin{remark}
Given the upper bound $ \edTime{n} $, we can convert the algorithm of \cref{lem:explicit.consensus.clustering} into one whose amortized update time satisfies \eqref{eq:explicit.consensus.time} on every execution, preserving the approximation guarantee with probability at least $ 1-\zeta $.
\end{remark}

\begin{definition}[Database model~\cite{assadi-wang}]\label{def:query.model.label}
At time $ t $, the algorithm accesses the graph of positive edges on $ V_t $ through three queries.
A degree query returns the number of neighbors of a given vertex, an edge query tests whether two given vertices are adjacent, and a neighborhood-sample query returns an independent uniformly random neighbor of a given vertex of positive degree.
Each query costs $ O(\log(|V_t|+2)) $ time.
\end{definition}

\paragraph{From clique buckets to the database model.}\label{par:explicit.bucket.queries}
By the bucket representation of \cref{subsec:offline-buckets,subsec:dynamic-pivot}, two distinct vertices of a three-input majority graph are adjacent exactly when their input labels agree on some pair of coordinates.
For a vertex $ v $, the intersection of any two of its three pair buckets is the set of vertices agreeing with $ v $ on all three coordinates.
Order each pair bucket by the remaining coordinate and then by vertex index in an order-statistic tree, so this intersection occupies one rank interval.
The open neighborhood is the disjoint union of the common intersection with $ v $ removed and the complement of the intersection in each of the three buckets.
Summing these four sizes gives the degree, and a uniform rank in their concatenation gives a uniform neighbor.
Thus insertions and all three database queries take $ O(\log(|V_t|+2)) $ time, using three bucket memberships per vertex.

\begin{restatable}{lemma}{ExplicitDA}\label{lem:explicit.correlation.clustering}
For some absolute constant $ \edFactor\ge1 $, the randomized algorithm $ \explicitDA $ in the database model~(\cref{def:query.model.label}) achieves an $ \edFactor $-approximation simultaneously at all prefixes of any fixed insertion-only stream with probability at least $ 1-N^{-2} $, where $ N=\max\{n,n_0\} $ for a sufficiently large absolute constant $ n_0 $.
%The algorithm uses $ N $ as its size parameter and takes the majority of three independent runs of each agreement or heaviness test.
For every $ t\le n $, the expected total work through insertion $ t $, including all reported label assignments, is $ O(t\polylog(n+2)) $.
\end{restatable}

\begin{proof}[Proof of~\cref{thm:explicit.dynamic}]
Apply \cref{lem:explicit.consensus.clustering} with the algorithm of \cref{lem:explicit.correlation.clustering}.
For fixed $ \eps $ and inverse-polynomial $ \zeta $, \eqref{eq:explicit.consensus.time} is polylogarithmic in $ n $.
\end{proof}

We defer the proof of \cref{lem:explicit.correlation.clustering} to \cref{sec:explicit_dynamic_correlation_clustering_algorithm_for_insertion_only_stream}.
The remainder of this section is devoted to prove \cref{lem:explicit.consensus.clustering}.
The cases $ n=1 $ and $ m\le2 $ are handled by following an input, as in \cref{sec:online}, so assume $ n\ge2 $ and $ m\ge3 $.

\subsection{Algorithm}\label{subsec:explicit.consensus.algorithm}

We follow the history-score selection rule in \cref{sec:oncon-algorithm}, with the following two modifications.
For every triple candidate, maintain its solution by employing $ \edSolver $ instead of $ \Agree $.
In \cref{sec:oncon-algorithm}, the algorithm decides whether to retain the selected candidate or switch to another candidate using history scores computed from exact costs.
As maintaining the exact cost of each candidate is too expensive in dynamic model, we maintain an estimated cost of each candidate using a set of evaluation samples. 
We detail this estimation below.

We reuse the parameter choices from \cref{sec:oncon-algorithm,subsec:oncon-analysis}:
\begin{equation*}
  c=\eta+1,\qquad
  w=\frac{3c}{3c+2},\qquad
  R_0=2-\frac{2}{3\eta+5},\qquad
  R=R_0+\frac{\eps}{2},
\end{equation*}
with $ \eta $ the approximation factor of $ \edSolver $. 
The candidate set is defined as in \cref{sec:oncon-algorithm}, with the parameters changed by a constant factor as follows.

\begin{align}
    \theta=\gamma=\frac{\eps}{40}, \qquad
    K=\left\lceil\frac{8}{\eps}\log\frac{4n}{\zeta}\right\rceil \label{eq:explicit.dynamic.consensus.K}
\end{align}

Sample $ K $ candidates independently as in \cref{sec:oncon-algorithm}: with probability $ w $, follow a uniform input, and otherwise run $ \edSolver $ on the majority graph of a uniform triple of distinct input indices.
Let $ \A=[K] $, and let $ \A_{\rm in} $ and $ \A_{\rm tri} $ be the sets of indices of input and triple candidates, respectively. As in \cref{sec:oncon-algorithm}, write $ A_a(t) $ for candidate $ a $'s partition and $ \lambda_t^a(v) $ for its local identifier of vertex $ v $.

\paragraph{Evaluation samples.}
Set
\begin{align}
    L=\left\lceil\log_2\frac{4n}{\zeta}\right\rceil, \qquad
    g=\left\lceil\frac{2}{\gamma^2}\log\frac{8KnL}{\zeta}\right\rceil \qquad
    s=Lg.
    \label{eq:explicit.consensus.parameters}
\end{align}

Independently of the candidates, we sample $ L $ groups $ (J_{j,1},\ldots,J_{j,g}) $, with all $ J_{j,b} $ independent and uniform in $ [m] $ as evaluation sets use throughout the course of the algorithm.
For a partition $ C $ of $ V_t $, define
\begin{align}
    \edmean{j}{t}(C)&:=\frac1g\sum_{b=1}^g d(C,P_{J_{j,b}}^t),
    &\edcost{t}(C)&:=\min_{j\in[L]}\edmean{j}{t}(C),\label{eq:explicit.consensus.estimates}\\
    \edscore{a}{t}&:=\edcost{t}(A_a(t)),
    &\edhistory{a}{t}&:=\max_{u\le t}\edscore{a}{u},
    &\edminimum{t}&:=\min_{a\in\A}\edhistory{a}{t}.\nonumber
\end{align}
Initially, $ \edhistory{a}{0}=0 $.

The updating rule of the algorithm are given in \cref{alg:explicit.consensus}.

\tienlong{Unreadable beyond this point. The details are correct though.}

\begin{algorithm}[htbp]
    \DontPrintSemicolon
    \KwIn{The arriving vertex $ v_t $, the sampled indices, and the maintained candidate and selection states.}
    \KwOut{The identifier of $ v_t $ and every changed output identifier.}
    Read and store the labels of $ v_t $ at the sampled input indices\;
    \ForEach{$ a\in\A $}{
        Advance candidate $ a $ to time $ t $; for a triple candidate, first update its majority-graph representation\;
        Let $ \edlog{a}{t} $ contain the assignment to $ v_t $ and the reported changes to old vertex labels\;
        Update all group costs by processing $ \edlog{a}{t} $\;
        Compute $ \edscore{a}{t} $ and set $ \edhistory{a}{t}\gets\max\{\edhistory{a}{t-1},\edscore{a}{t}\} $\;
    }
    Set $ \edminimum{t}\gets\min_{a\in\A}\edhistory{a}{t} $ and apply the selection rule in \cref{sec:oncon-algorithm}, using estimated history scores, to choose $ a_t $\;
    \eIf{$ t=1 $ or $ a_t\ne a_{t-1} $}{
        Assign and report $ (a_t,\lambda_t^{a_t}(v)) $ for every $ v\in V_t $\;
    }{
        Apply and report the assignments in $ \edlog{a_t}{t} $, prefixing each local label by $ a_t $\;
    }
    \caption{Explicit dynamic consensus-clustering algorithm.}
    \label{alg:explicit.consensus}
\end{algorithm}

\subsection{Approximation Guarantee}\label{subsec:explicit.consensus.approximation}

For the analysis, define
\begin{align}
    \edproxy{a}{t}:=
    \begin{cases}
        \cost_t(P_{i_a}^t),&a\in\A_{\rm in},\\
        \OPT_t+c|B_{Q_a}^t(F_t^*)|,&a\in\A_{\rm tri}.
    \end{cases}
    \label{eq:explicit.consensus.proxy}
\end{align}
Call $ a $ good for $ t $ if $ \edproxy{a}{t}\le R\OPT_t $ and, when $ a $ is a triple candidate, its solver achieves the $ \edFactor $-approximation at every prefix.

\begin{lemma}[Existence of a good candidate]\label{lem:explicit.good.candidate}
With probability at least $ 1-\zeta/4 $, every prefix $ t\le n $ has a good candidate.
\end{lemma}
\begin{proof}
The proxy depends only on the sampled candidate type and input indices.
The input expectation in \eqref{eq:oncon-input-history}, the majority-error bound in \eqref{eq:oncon-majority-error}, and the cancellation in the proof of \cref{lem:oncon-expectation} give $ \E[\edproxy{a}{t}]\le R_0\OPT_t $.
The fixed-prefix Markov argument in the proof of \cref{lem:oncon-good-history} therefore gives
\begin{align}
    \Prb[\edproxy{a}{t}\le R\OPT_t]\ge\frac{\eps}{4}.
    \label{eq:explicit.consensus.proxy.success}
\end{align}
If $ \OPT_t=0 $, the proxy is zero for every draw, so this bound still holds.

Conditional on any fixed triple, its solver succeeds throughout the stream with probability at least $ 1/2 $.
Consequently, a candidate is good for $ t $ with probability at least $ \eps/8 $, independently across candidates.
By the choice of $ K $ in \eqref{eq:explicit.dynamic.consensus.K}, the probability that any prefix has no good candidate is at most
\begin{align}
    n\exp(-\eps K/8)\le\zeta/4.
    \label{eq:explicit.consensus.candidate.event}
\end{align}
\end{proof}

\begin{lemma}[Accuracy of the estimated costs]\label{lem:explicit.estimated.costs}
With probability at least $ 1-\zeta/2 $, simultaneously for every $ t\le n $ and $ a\in\A $,
\begin{align}
    \edcost{t}(F_t^*)&\le2\OPT_t,\label{eq:explicit.consensus.ceiling}\\
    \left|
        \edcost{t}(A_a(t))-\edcost{t}(F_t^*)
        -\bigl(\cost_t(A_a(t))-\OPT_t\bigr)
    \right|
    &\le\gamma d(A_a(t),F_t^*).
    \label{eq:explicit.consensus.comparison}
\end{align}
\end{lemma}
\begin{proof}
To prove \eqref{eq:explicit.consensus.ceiling}, apply Markov's inequality to each group mean for the fixed partition $ F_t^* $.
Each mean exceeds $ 2\OPT_t $ with probability at most $ 1/2 $, so their minimum does so with probability at most $ 2^{-L} $.
When $ \OPT_t=0 $, all these means are zero.
A union bound over $ t $ gives failure probability at most $ \zeta/4 $.

For \eqref{eq:explicit.consensus.comparison}, condition on all candidate samples and solver choices, leaving the evaluation samples independent.
For each $ a,t,j $, the difference of the two group means averages independent variables $ d(A_a(t),P_{J_{j,b}}^t)-d(F_t^*,P_{J_{j,b}}^t) $.
By the same triangle-inequality argument used in \cref{lem:amplification-selection}, each variable lies in $ [-d(A_a(t),F_t^*),d(A_a(t),F_t^*)] $.
Hoeffding's inequality bounds the probability of a deviation larger than $ \gamma d(A_a(t),F_t^*) $ in either direction by $ 2\exp(-g\gamma^2/2) $.
If the distance is zero, the comparison is exact.
A union bound over $ KnL $ comparisons gives failure probability at most $ \zeta/4 $.
Finally, if the differences of two lists of numbers all lie in the same interval, the difference of their minima lies in that interval as well.
Taking the minimum over the groups proves \eqref{eq:explicit.consensus.comparison}.

\end{proof}

\begin{lemma}[Estimated history of a good candidate]\label{lem:explicit.estimated.history}
If $ b $ is good for $ t $ and \eqref{eq:explicit.consensus.comparison} holds at $ t $, then
\begin{align}
    \edminimum{t}\le\edhistory{b}{t}
    \le\edcost{t}(F_t^*)+(R-1+3\gamma)\OPT_t.
    \label{eq:explicit.consensus.good.history}
\end{align}
\end{lemma}
\begin{proof}
The first inequality follows from the definition of $ \edminimum{t} $.
If $ b $ is an input candidate, each group cost is nondecreasing under insertions, so $ \edhistory{b}{t}=\edcost{t}(P_{i_b}^t) $ by the input-cost monotonicity argument in the proof of \cref{lem:oncon-expectation}.
Since $ \cost_t(P_{i_b}^t)\le R\OPT_t $, the triangle inequality gives $ d(P_{i_b}^t,F_t^*)\le(R+1)\OPT_t\le3\OPT_t $.
Substituting these bounds in \eqref{eq:explicit.consensus.comparison} proves the second inequality for an input candidate.

For a successful triple candidate, the derivation of \eqref{eq:oncon-history-distance} uses only the solver's simultaneous approximation guarantee, so it gives $ d(A_b(u),F_t^*|_{V_u})\le c|B_{Q_b}^t(F_t^*)| $ for every $ u\le t $.
Each group mean, and hence their minimum, satisfies both the cost comparison and restriction inequalities used in \cref{lem:oncon-history}.
Applying that proof with estimated costs yields
\begin{align}
    \edhistory{b}{t}
    \le\edcost{t}(F_t^*)+c|B_{Q_b}^t(F_t^*)|
    \le\edcost{t}(F_t^*)+(R-1)\OPT_t.
    \nonumber
\end{align}
\end{proof}

\begin{lemma}[Approximation at every prefix]\label{lem:explicit.selected.cost}
With probability at least $ 1-\zeta $, the output satisfies $ \cost_t(C_t)\le(R_0+\eps)\OPT_t $ simultaneously for every $ t\le n $.
\end{lemma}
\begin{proof}
By \cref{lem:explicit.good.candidate,lem:explicit.estimated.costs} and a union bound, with probability at least $ 1-3\zeta/4 $ every prefix has a good candidate and satisfies \eqref{eq:explicit.consensus.ceiling} and \eqref{eq:explicit.consensus.comparison}.
On this event, fix $ t $ and apply \cref{lem:explicit.estimated.history} to a good candidate at $ t $.
The proof of \cref{lem:oncon-selection} uses only the monotonicity of history scores and the fact that a current score is at most its history score.
It therefore applies without change and gives $ \edcost{t}(C_t)\le(1+\theta)\edminimum{t} $.
Combine this inequality with \eqref{eq:explicit.consensus.good.history}, \eqref{eq:explicit.consensus.comparison} for the selected candidate, and $ d(C_t,F_t^*)\le\cost_t(C_t)+\OPT_t $ to obtain
\begin{align}
    (1-\gamma)\cost_t(C_t)
    &\le(1+\gamma)\OPT_t+\theta\edcost{t}(F_t^*)
        +(1+\theta)(R-1+3\gamma)\OPT_t\nonumber\\
    &\le\bigl(R+4\gamma+\theta(R+1+3\gamma)\bigr)\OPT_t.
    \label{eq:explicit.consensus.final.cost}
\end{align}
The last inequality uses \eqref{eq:explicit.consensus.ceiling}.
Since $ \theta=\gamma=\eps/40<1/160 $ and $ R<2 $,
\begin{align}
    \frac{R+4\gamma+\gamma(R+1+3\gamma)}{1-\gamma}
    &=R+\frac{\gamma(2R+5+3\gamma)}{1-\gamma}
    <R+10\gamma
    =R_0+\frac{3\eps}{4}
    <R_0+\eps.
    \nonumber
\end{align}
This holds for every prefix on the stated event, including prefixes with $ \OPT_t=0 $.
\end{proof}

\subsection{Implementation and Update Time}\label{subsec:explicit.consensus.implementation}


\begin{lemma}[Maintaining the estimated costs]\label{lem:explicit.cost.maintenance}
For each candidate, its $ L $ evaluation-group means can be maintained in $ O(ns) $ words and $ O(s\log(n+2)) $ time per reported label assignment, where $ s=Lg $.
Computing the estimated cost from the updated means takes $ O(L) $ time.
\end{lemma}
\begin{proof}
Apply the counters and update formulas of \cref{subsec:dynamic-costs} separately to each evaluation group, with group size $ g $.
Restore the input-only term following \eqref{eq:dynamic-score} to obtain the full group mean before taking its minimum or updating a history score.
That term cannot be omitted here, since it varies across groups and times.
For each candidate, keep its previous labels while processing the arriving vertex's assignment and then the old vertices' reported changes.
A change to a previously unused label is an insertion into an empty block.
As in \cref{subsec:dynamic-costs}, each assignment or move takes $ O(s\log(n+2)) $ time, including updates to all group sums.
Each of the $ s $ evaluation positions has at most $ n $ nonempty block-intersection records, giving the stated space bound.
Computing the minimum over the $ L $ groups takes $ O(L) $ time per candidate per arrival.
\end{proof}

\begin{lemma}[Cost of rebuilding the explicit output]\label{lem:explicit.rebuilding}
On every execution, the number $ S $ of outer switches satisfies
\begin{align}
    S\le K\left(2+\left\lceil\log_{1+\theta}(gn^2)\right\rceil\right)
    =O\!\left(K\eps^{-1}\log(gn+2)\right).
    \label{eq:explicit.consensus.switches}
\end{align}
Over the first $ t $ insertions, rebuilding the output at switches takes $ O(t(S+1)) $ time.
\end{lemma}
\begin{proof}
Every estimated score, history score, and minimum history score belongs to $ g^{-1}\mathbb Z\cap[0,\binom n2] $.
Applying the phase-count argument in the proof of \cref{lem:oncon-recourse}, with minimum positive score $ 1/g $ and the present value of $ \theta $, gives \eqref{eq:explicit.consensus.switches}.
Each rebuild writes and reports at most $ t $ identifiers during the first $ t $ insertions.
Between switches, output writes and reports are charged to assignments reported by the selected candidate.
\end{proof}

\begin{lemma}[Total update time]\label{lem:explicit.total.work}
For every $ t\le n $, the expected total time for the first $ t $ insertions, including initialization and all output reports, is
\begin{align}
    O\!\left(
        tK(\edTime{n}+1)s\log(n+2)
        +tK\eps^{-1}\log(gn+2)
    \right).
    \label{eq:explicit.consensus.total.work}
\end{align}
\end{lemma}
\begin{proof}
For every fixed choice of candidate samples, the expected work of all solver copies over the first $ t $ insertions, including their queries and reported assignments, is $ O(tK\edTime{n}) $.
The bucket representation implements each database query in $ O(\log(n+2)) $ time, and \cref{lem:explicit.cost.maintenance} charges $ O(s\log(n+2)) $ time to each reported assignment.
Reading the sampled labels, updating input candidates, and computing scores and selections fit in $ O(tKs\log(n+2)) $ time.
By \cref{lem:explicit.rebuilding}, output rebuilds contribute $ O(tK\eps^{-1}\log(gn+2)) $ time.
Adding these bounds proves \eqref{eq:explicit.consensus.total.work}.
\end{proof}

\begin{lemma}[Bounding the work of each solver copy]\label{lem:explicit.capped.work}
Given the prefix expectation bound in \cref{lem:explicit.consensus.clustering}, one can modify \cref{alg:explicit.consensus} so that \eqref{eq:explicit.consensus.time} bounds the amortized update time at every prefix on every execution.
The modified algorithm retains the simultaneous approximation guarantee of \cref{lem:explicit.selected.cost}.
\end{lemma}
\begin{proof}
Double $ K $ in \eqref{eq:explicit.dynamic.consensus.K} and use this larger value when setting $ g $.
Set $ h=1+\lceil\log_2 n\rceil $.
At insertion $ t $, stop a triple solver permanently if its next operation would bring its cumulative work above $ 8ht\edTime{n} $.
Count only the solver's own work against its budget, and buffer its reports until the insertion finishes.
If the solver stops, discard its pending reports and use the first input partition for that candidate from this insertion onward.
Cache that input's labels as vertices arrive, and rebuild the candidate's labels and score counters once when it switches to that input.
Use fresh identifiers for its input blocks, report the resulting assignments, and retain the candidate's previous history score when taking the running maximum.

Fix a sampled triple and let $ W(u) $ be the cumulative work of its uninterrupted solver through insertion $ u $.
For any $ t\le n $, the first endpoint $ u\ge t $ in $ \{\min(2^j,n):0\le j<h\} $ satisfies $ u\le2t $.
If the budget is exceeded at $ t $, monotonicity gives $ W(u)>4hu\edTime{n} $.
Since $ \E[W(u)]\le u\edTime{n} $, Markov's inequality and a union bound over the at most $ h $ endpoints bound the probability of stopping by $ 1/4 $.
Thus, conditional on any triple, its solver both succeeds at all prefixes and finishes the entire stream with probability at least $ 1/2-1/4=1/4 $.
Using this probability in the proof of \cref{lem:explicit.good.candidate} gives success probability at least $ \eps/16 $ per candidate at each fixed prefix, for which the doubled $ K $ suffices.
The modified trajectories are independent of the evaluation samples, so \cref{lem:explicit.estimated.costs,lem:explicit.estimated.history,lem:explicit.selected.cost} apply as before.

Through insertion $ t $, each solver spends at most $ 8ht\edTime{n} $ operations and rebuilds at most once, at a cost of $ O(ts\log(n+2)) $ for that rebuild.
Consequently the accounting of \cref{lem:explicit.total.work} applies on every execution with $ 8h\edTime{n} $ in place of $ \edTime{n} $.
The additional factor $ h=O(\log(n+2)) $ is absorbed by the polylogarithmic factor in \eqref{eq:explicit.consensus.time}.
\end{proof}

\begin{proof}[Completion of the proof of~\cref{lem:explicit.consensus.clustering}]
The approximation guarantee follows from \cref{lem:explicit.selected.cost}.
By \eqref{eq:explicit.dynamic.consensus.K} and \eqref{eq:explicit.consensus.parameters},
\begin{align}
    K&=O\!\left(\eps^{-1}\log\frac{2n}{\zeta}\right),
    &L&=O\!\left(\log\frac{2n}{\zeta}\right),
    &g&=O\!\left(\eps^{-2}\log\frac{2n}{\eps\zeta}\right).
    \nonumber
\end{align}
Substituting $ s=Lg $ and $ \edTime{n}=\polylog(n+2) $ into \eqref{eq:explicit.consensus.total.work} proves \eqref{eq:explicit.consensus.time}.
\end{proof}

\section{Explicit Dynamic Correlation Clustering Algorithm for Insertion-Only Stream}\label{sec:explicit_dynamic_correlation_clustering_algorithm_for_insertion_only_stream} % (fold)

In this section, we prove the following lemma.
\ExplicitDA*

\tienlong{Note to myself: for whatever changes and fixes we made to DA, make sure we mention that the other part are unaffected}
\subsection{Algorithm Description}\label{sub:algorithm_description} % (fold)


Dynamic Agreement Algorithm that reports explicit updated vertices.

At the end of the Dynamic Agreement Algorithm, we call the LabelSync algorithm to report the updated vertices and their new labels.

To support that, we use the following additional data structures:
\begin{itemize}
    \item A journal $ \journal $ records every edge insertion, edge deletion, and change of anchor membership during update $ t $.
    For each edge modified during the update, the journal also records whether it was present immediately before its first modification.
    \tienlong{The journal should be a data structrure of the $ \explicitDA $, not $ \LabelSync $}
\item We use the dynamic connectivity structure of~\cite{holm2001poly} to maintain the connected components of the sparse graph $ \sparseG $ under edge insertions and deletions. This structure supports a query $ \comp{x} $ in $ O(\log n) $ time with the property that $ \comp{x}=\comp{y} $ if and only if $ x $ and $ y $ belong to the same component of the current graph. We refer to $ \comp{x} $ as the \emph{component identifier} of $ x $.
\end{itemize}

\tienlong{Remember to add input and output}
\begin{algorithm}[htbp]
    \SetKwInOut{KwIn}{Input}
    \SetKwInOut{KwOut}{Output}
    \KwIn{}
    \KwOut{}
        A maintained sparsified graph $ \sparseG \gets \emptyset $

        A data structure to maintain connected components of $ \sparseG $. \tcc{{\color{red} Citation}}

        \New{A journal to record the updates to $ \sparseG $: $ \journal \gets \emptyset $}

        \On{arrival of $v_t$}{
            $ \Notify{v_{t}, \epsDA } $

            \ForAll{$ u \in \interestSet{t} $}{
                $ \Clean{ u, \epsDA, t } $

                $ \Anchor{ u, \epsDA, t } $

                $ \Connect( u, \epsDA, t ) $
            }

        }

        \New{$ L_{t} \gets \{ u \in \interestSet{t} \cap \anchorSet:\ h_{u} = \boolNo, \text{ or } \card{\closedN{\sparseG}{u}} \leq (1-\epsDA)\card{\closedN{\localG{t}}{u}} \} $} \tcc{Note that this Check for closed neighborhood}
        
        \New{delete every vertex of $ L_{t} $ from $ \anchorSet $}

        \New{delete every vertex of $ \sparseG $ with no endpoint in $ \anchorSet $}

        \New{
            \ForAll{$ u \in \interestSet{t} \cap \anchorSet $}{
                $ \snapshotN{u}\gets \openN{\sparseG}{u} $
            }
        }
        \New{$ \LabelSync{t} $}

        \KwRet{}
    \caption{$ \explicitDA $: Explicit Dynamic Agreement Algorithm for Insertion-Only Stream}
    \label{alg.explicit.dynamic.agreement}
\end{algorithm}

\paragraph{\texorpdfstring{$ \LabelSync $}{LabelSync} Algorithm}\label{par:_texorpdfstring_labelsync_labelsync_algorithm} % (fold)
At the end of update $ t $, $ \LabelSync $ maintains a partition $ \ccclasses{t} $ of the current vertex set $ V_t $.
We call each member of $ \ccclasses{t} $ a \emph{class}.
Each class $ \ccclass $ has a distinct label $ \ell(\ccclass) $, and every vertex $ w\in\ccclass $ stores this label as $ \ell_t(w)=\ell(\ccclass) $.
The partition $ \ccclasses{t} $ is the solution reported by $ \explicitDA $.
At the end of each update $ t $, the algorithm reports only the vertices that changed clusters, that is, $ (v_t,\ell_t(v_t)) $ and every $ (w,\ell_t(w)) $ with $ w\in V_{t-1} $ and $ \ell_t(w)\ne\ell_{t-1}(w) $.
We prove in \cref{sub:analysis_for_labelsync} that, on the success event $ \successDA $, the classes in $ \ccclasses{t} $ are exactly the connected components of $ \sparseG{t} $.
% paragraph LabelSync Algorithm (end)

\paragraph{Data structures for \texorpdfstring{$ \LabelSync $}{LabelSync}}\label{par:data_structure_of_labelsync}
To implement $ \LabelSync $, we maintain the following data structures.
\begin{itemize}
    \item For each class $ \ccclass $, we store its size and a doubly linked list of its vertices.
    \item For each vertex $ w $ of a class $ \ccclass $, we stores its label and a pointer to its position in the doubly linked list of $ \ccclass $. Additionaly, $ w $ stores its membership in $ \anchorSet{t} $ and its degree $ \deg_t(w):=\card{\openN{\sparseG{t}}{w}} $ in the sparse graph.
    These values, together with the class lists and vertex labels, are updated only in the final step of $ \LabelSync $.
    Until that step, they refer to time $ t-1 $.
\end{itemize}

\paragraph{Description of \texorpdfstring{$ \LabelSync $}{LabelSync} Algorithm}\label{par:description_of_texorpdfstring_labelsync_labelsync_algorithm} % (fold)
The input to $ \LabelSync $ includes the sparse graph $ \sparseG{t} $ (after the update from time $ t-1 $ to time $ t $), the journal $ \journal $ of update $ t $, the arriving vertex $ v_t $, the data structure of $ \LabelSync $ at time $ t-1 $.

The goal is to update the data structure to reflect the connected components of $ \sparseG{t} $ and to report the vertices whose labels changed.

On a high level, the algorithm achieves this by performing a \emph{split} and a \emph{merge} operation. Suppose that at time step $ t-1 $, $ \LabelSync $ successfully maintained the $ \ccclasses{t-1} $ as the connected components of $ \sparseG{t-1} $. At time step $ t $, vertices in some components in $ \ccclasses{t-1} $ may not belong to the same components in $ \sparseG{t} $, and some vertices in different components in $ \ccclasses{t-1} $ may belong to the same component in $ \sparseG{t} $. To fix this, we first perform the split operation, which partitions each old class $ \ccclass \in \ccclasses{t-1} $ into \piece{s}, where each \piece{} is a subset of vertices of $ \ccclass $ that belong to the same component in $ \sparseG{t} $. After that, the merge operation groups \piece{s} that belong to the same component in $ \sparseG{t} $ together. Finally, we update the labels of vertices in the new classes and report the vertices whose labels changed.

\tienlong{May be briefly mention we need to design the algorithm so that it operates mostly on the updated vertices and edges.}
In more details, the algorithm consists of the four following steps.

\noindent\textbf{Step 1.}
From the journal, we compute
\begin{align}
    \delt{t}:=E(\sparseG{t-1})\setminus E(\sparseG{t}), \nonumber
    \qquad
    \addt{t}:=E(\sparseG{t})\setminus E(\sparseG{t-1}), \nonumber
\end{align}
the sets of edges deleted and added during update $ t $, respectively.
These sets can be computed in one pass over the journal.
For each old class $ \ccclass\in\ccclasses{t-1} $, define
\[
    D_{\ccclass}:=\{(x,y) \in\delt{t}:x,y\in\ccclass\},
    \qquad d_{\ccclass}:=\card{D_{\ccclass}}.
\]
For each $ w\in\ccclass $, define
\[
    \delIncEdge{w}:=\card{\{e\in D_{\ccclass}:w\in e\}}.
\]
We can construct the lists $ D_{\ccclass} $ and accumulate the counts $ d_{\ccclass} $ and $ \delIncEdge{w} $ in one pass over $ \delt{t} $. \tienlong{Probably mention when analyze the running time that this step can be implemented using hashmap}

\noindent\textbf{Step 2 (Split).}
We partition each old class into temporary sets called \emph{\piece{s}}, such that each \piece{} lies entirely within one component of $ \sparseG{t} $.

To do that, observe that conditioned on $ \successDA $, the \piece{s} of an old class $ \ccclass $ will be exactly
\begin{align}
    \mathcal P_{\ccclass}
    =\{\ccclass\cap K:K\text{ is a component of }\sparseG{t},\ \ccclass\cap K\ne\emptyset\}. \nonumber
\end{align}
Among \piece{s} in $ \mathcal P_{\ccclass} $, we choose one \piece{} to inherit the old label $ \ell(\ccclass) $. We refer to this \piece{} as the \emph{retained \piece{}} of $ \ccclass $, and denote it by $ \retainedPieces{\ccclass} $.
The other \piece{s} initially have no assigned label.

If $ d_{\ccclass}=0 $, the class is represented by the single retained \piece{} $ R_{\ccclass}=\ccclass $, without listing its vertices.
For each class with $ d_{\ccclass}>0 $, we distinguish two cases.
\begin{itemize}
    \item If $ d_{\ccclass}\ge\card{\ccclass}/2-1 $, we traverse $ \ccclass $ and group its vertices into \piece{s} such that, $ x,y \in \ccclass $ are grouped into the same \piece{} if and only if $ \comp{x}=\comp{y} $.

    A largest \piece{}  is chosen as $ \retainedPieces{\ccclass} $.
    % Since $ \card{\ccclass}\le 2d_{\ccclass}+2\le 4d_{\ccclass} $, the number of vertices examined is proportional to $ d_{\ccclass} $.
    \item If $ d_{\ccclass}<\card{\ccclass}/2-1 $, choose an endpoint $ r\in\ccclass\cap\anchorSet{t-1} $ of an edge in $ D_{\ccclass} $.
    Such an endpoint exists because every old edge has an endpoint in $ \anchorSet{t-1} $, by definition of the anchor set.
    Define
    \begin{align}
        W_{\ccclass}:=\{w\in\bigcup_{e\in D_{\ccclass}}e:
        w \in \ccclass\setminus\anchorSet{t-1},\ \delIncEdge{w}=\deg_{t-1}(w)\}, \nonumber
    \end{align}
    the set of vertices in $ \ccclass\setminus \anchorSet{t-1} $ whose all edges edges in time $ t-1 $ are deleted in time $ t $.
    We will show that (\cref{lem:labelsync.robust.components}), on $ \successDA $, every vertex of $ \ccclass\setminus W_{\ccclass} $ remains connected to $ r $ after the deletions.
    For each vertex $ w\in W_{\ccclass} $, we group it with the same \piece{} as $ r $ if and only if $ \comp{w}=\comp{r} $.
    For the remaining vertices in $ W_{\ccclass} $, we group them into \piece{s}
    \begin{align}
        P_{\ccclass,\kappa}:=\{w\in W_{\ccclass}:\comp{w}=\kappa\}, \nonumber
    \end{align}
    for each $ \kappa\ne\comp{r} $ attained in $ W_{\ccclass} $.
    The retained \piece{} is
    \[
        R_{\ccclass}:=\ccclass\setminus\bigcup_{\kappa\ne\comp{r}}P_{\ccclass,\kappa},
    \]
    where the union ranges only over the identifiers attained in $ W_{\ccclass} $.
    % We store $ R_{\ccclass} $ implicitly, with representative $ r $ and size $ \card{\ccclass}-\sum_{\kappa\ne\comp{r}}\card{P_{\ccclass,\kappa}} $.
    % The set $ W_{\ccclass} $ is found from the endpoints of $ D_{\ccclass} $, and its vertices are grouped using a hash map indexed by their component identifiers.
    % Thus this case does not require listing $ \ccclass $.
    We note that this case can be implemented by traversing $ D_{\ccclass} $ and $ W_{\ccclass} $, instead of traversing the entire class $ \ccclass $.
\end{itemize}

The \piece{s} of all old classes, together with the singleton $ \{v_t\} $, form a partition $ \mathcal P $ of $ V_t $.
% The singleton $ \{v_t\} $ initially has no assigned label.
% For each represented \piece{} $ P $, we store its size, a representative $ x_P\in P $, and its inherited label $ \lambda(P) $, where $ \lambda(P)=\bot $ if it has none.
% The inherited labels are temporary attributes of \piece{s}; vertex labels remain unchanged until Step 4.
% Parts other than implicitly represented retained \piece{s} are stored as explicit vertex lists.
% A hash map records the \piece{} containing each old vertex outside a retained part.
% Every other old vertex belongs to the retained \piece{} of its old class, so the \piece{} containing any given vertex can be found without listing its class.

\tienlong{Auditing has not been performed beyond this point.}
\noindent\textbf{Step 3 (Merge).}
Let $ \mathcal Z $ contain all \piece{s} constructed for classes with $ d_{\ccclass}>0 $, the singleton $ \{v_t\} $, and the \piece{} containing each endpoint of an edge in $ \addt{t} $.
Each \piece{} occurs only once in $ \mathcal Z $.
If an endpoint belongs to a class with $ d_{\ccclass}=0 $, we create its \piece{} record when needed, using that endpoint as its representative and $ \ell(\ccclass) $ as its inherited label.
For each $ P\in\mathcal Z $, define its \emph{component key} by $ \mathrm{key}(P):=\comp{x_P} $.
We group the \piece{s} by their component keys using a hash map.
For each key $ \kappa $ that occurs, choose a largest \piece{} $ P^* $ with that key, breaking ties arbitrarily.
If $ \lambda(P^*)\ne\bot $, all \piece{s} with key $ \kappa $ are assigned $ \lambda(P^*) $.
Otherwise, they are assigned a previously unused label.
Classes whose \piece{s} are outside $ \mathcal Z $ keep their old labels.

\noindent\textbf{Step 4 (Report).}
For every \piece{} whose assigned label differs from the old label of its vertices, we list its vertices and collect their new labels.
We also collect the label assigned to $ v_t $.
An implicit retained \piece{} is listed by traversing its old class and omitting vertices recorded in other \piece{s} of that class.
After all such traversals, we apply the collected changes and update the class lists and sizes.
We then refresh the stored degrees and anchor memberships at vertices appearing in $ \journal $ and at $ v_t $; all other stored values remain valid.
Finally, we report $ (v_t,\ell_t(v_t)) $, followed by every $ (w,\ell_t(w)) $ with $ w\in V_{t-1} $ and $ \ell_t(w)\ne\ell_{t-1}(w) $.
% paragraph Description of LabelSync Algorithm (end)

\begin{algorithm}[htbp]
  \small
  \DontPrintSemicolon
  \SetKwInOut{KwIn}{Input}
  \SetKwInOut{KwOut}{Output}
  \KwIn{The journal $ \journal $ of update $ t $, the arriving vertex $ v_t $, the stored state at time $ t-1 $, and connectivity queries in $ \sparseG{t} $.}
  \KwOut{$ (v_t,\ell_t(v_t)) $ and every $ (w,\ell_t(w)) $ with $ w\in V_{t-1} $ and $ \ell_t(w)\ne\ell_{t-1}(w) $.}

  Compute $ \delt{t} $, $ \addt{t} $, the lists $ D_{\ccclass} $, and the counters $ d_{\ccclass} $ and $ \delIncEdge{w} $ from $ \journal $\;
  $ \piecesToMerge\gets\{\{v_t\}\} $, with $ \lambda(\{v_t\})=\bot $\;

  \tcc{Split old classes; vertex labels are not changed.}
  \ForEach{$ \ccclass $ with $ d_{\ccclass}>0 $}{\nllabel{line:labelsync.split.begin}
    \uIf{$ d_{\ccclass}\ge \card{\ccclass}/2-1 $}{\nllabel{line:labelsync.split.large.deletion.begin}
      $ \mathcal K_{\ccclass}\gets\{\comp{x}:x\in \ccclass\} $\;
      $ P_{\ccclass,\kappa}\gets\{x\in \ccclass:\comp{x}=\kappa\} $ for each $ \kappa\in\mathcal K_{\ccclass} $\;
      $ \mathcal P_{\ccclass}\gets\{P_{\ccclass,\kappa}:\kappa\in\mathcal K_{\ccclass}\} $\;
      Choose $ R_{\ccclass}\in\arg\max_{P\in\mathcal P_{\ccclass}}|P| $\;
      Store the \piece{s} as explicit lists, each with a representative\;
      \nllabel{line:labelsync.split.large.deletion.end}
    }
    \Else{\nllabel{line:labelsync.split.few.deletion.begin}
      Choose $ r\in \ccclass\cap\anchorSet{t-1} $ incident to an edge of $ D_{\ccclass} $\;
      $ W_{\ccclass}\gets\{w\in\bigcup_{e\in D_{\ccclass}}e: w\notin\anchorSet{t-1},\ \delIncEdge{w}=\deg_{t-1}(w)\} $\;
      $ \kappa_r\gets\comp{r} $\;
      $ \mathcal K_{\ccclass}\gets\{\comp{w}:w\in W_{\ccclass}\}\setminus\{\kappa_r\} $\;
      $ P_{\ccclass,\kappa}\gets\{w\in W_{\ccclass}:\comp{w}=\kappa\} $ for each $ \kappa\in\mathcal K_{\ccclass} $\;
      $ R_{\ccclass}\gets \ccclass\setminus\bigcup_{\kappa\in\mathcal K_{\ccclass}}P_{\ccclass,\kappa} $\;
      $ \mathcal P_{\ccclass}\gets\{R_{\ccclass}\}\cup \{P_{\ccclass,\kappa}:\kappa\in\mathcal K_{\ccclass}\} $\;
      Store each $ P_{\ccclass,\kappa} $ explicitly; represent $ R_{\ccclass} $ implicitly, with representative $ r $ and size $ \card{\ccclass}-\sum_{\kappa\in\mathcal K_{\ccclass}}|P_{\ccclass,\kappa}| $\;
    }
    $ \lambda(R_{\ccclass})\gets\ell(\ccclass) $; $ \lambda(P)\gets\bot $ for each $ P\in\mathcal P_{\ccclass}\setminus\{R_{\ccclass}\} $\;
    $ \piecesToMerge\gets\piecesToMerge\cup\mathcal P_{\ccclass} $\;
  }\nllabel{line:labelsync.split.end}

  \tcc{Merge \piece{s} with the same component key.}
  \tcp{May be rewrite the following line into:}
  \tcp{$ \piecesToMerge\gets\piecesToMerge\cup\{P: P \text{ is a \piece{} containing a vertex } x\in \addt{t} \text{ and } P \text{ is not already in }\piecesToMerge\} $, where $ P(x) $ is the \piece{} containing $ x $}
  \ForEach{endpoint $ x $ of an edge in $ \addt{t} $}{\nllabel{line:labelsync.merge.begin}
    $ \piecesToMerge\gets\piecesToMerge\cup\{P(x)\} $, where $ P(x) $ is the \piece{} containing $ x $\; \nllabel{line:labelsync.piecesToMerge.include.addt}
  }
  \ForEach{$ \pieceToMerge \in\piecesToMerge $}{
    $ \mathrm{key}(\pieceToMerge )\gets\comp{x_P} $\;\nllabel{line:labelsync.assign.key}
  }
  \ForEach{distinct key $ \kappa $ represented in $ \piecesToMerge $}{
    $ \piecesToMerge_\kappa\gets\{\pieceToMerge \in\piecesToMerge:\mathrm{key}(\pieceToMerge )=\kappa\} $\;\nllabel{line:labelsync.group.key}
    Choose $ \pieceToMerge^*\in\arg\max_{\pieceToMerge \in\piecesToMerge_\kappa} \card{\pieceToMerge} $\;
    $ L_\kappa\gets\lambda(\pieceToMerge ^*) $ if $ \lambda(\pieceToMerge ^*)\ne\bot $; otherwise let $ L_\kappa $ be a fresh label\;
    $ \widehat\ell(\pieceToMerge )\gets L_\kappa $ for every $ \pieceToMerge \in\piecesToMerge_\kappa $\;\nllabel{line:labelsync.assign.label}
  }\nllabel{line:labelsync.merge.end}

  \tcc{Apply the assignments and report changes.}
  Collect all required vertex-label changes before modifying any class list\;
  Apply the collected changes and update the class lists and sizes; classes outside $ \piecesToMerge $ keep their labels\;
  Refresh stored degrees and anchor flags only at vertices appearing in $ \journal $, and at $ v_t $\;
  \Return{$ (v_t,\ell_t(v_t)) $, followed by all $ (w,\ell_t(w)) $ with $ w\in V_{t-1} $ and $ \ell_t(w)\ne\ell_{t-1}(w) $}\;
  \caption{$ \LabelSync $ at the end of update $ t $.}
  \label{alg.label.sync}
\end{algorithm}


\subsection{Analysis for \texorpdfstring{$ \LabelSync $}{LabelSync}}
\label{sub:analysis_for_labelsync}

{\color{red} As shown in~\cite{CLMP2024Dynamic}}, the following invanriant holds for every time step $ t $:
\begin{align}
    \text{every edge of } \sparseG{t} \text{ has an endpoint in } \anchorSet{t}. \tag{I} \label{eq:invariant.anchor}
\end{align}

\paragraph{Correctness.}
\label{par:correctness}
We use the invariant that every edge of $ \sparseG{t} $ has an endpoint in $ \anchorSet{t} $.
Let $ \successDA $ denote the success event from the analysis of $ \explicitDA $; on this event, the density bound in \cref{lem:ccclass.is.dense} holds at every time step.

\begin{lemma} \label{lem:labelsync.robust.components}
Let $ \sparseG $ be a simple undirected graph and let $ \anchorSet $ be a vertex set containing at least one endpoint of every edge of $ \sparseG $.
Let $ \ccclass $ be a connected component with $ \card{\ccclass}>1 $, and suppose that
\[
  |\openN{\sparseG}{a}|\ge\tfrac34\card{\ccclass}-1
  \qquad\text{for every }a\in \ccclass\cap\anchorSet.
\]
For any $ D\subseteq E(\sparseG[\ccclass]) $ with $ |D|<\card{\ccclass}/2-1 $, the graph $ \sparseG[\ccclass]\setminus D $ has one component containing $ \ccclass\cap\anchorSet $.
Every other component is a singleton outside $ \anchorSet $, and there are at most $ |D| $ such singletons.
\end{lemma}
\begin{proof}
    As $ \card{\ccclass} > 1 $, $ \ccclass $ has an least one edge.
    By definition, $ \anchorSet $ contains at least one vertex of all edges.
    Hence, $ \ccclass \cap \anchorSet \neq \emptyset $.
    We proceed by cases.

    $ \bullet $ If $ \card{\ccclass \cap \anchorSet} = 1 $, suppose $ \ccclass \cap \anchorSet = \{ a \} $.
    By definition of $ \anchorSet $, this implies that all edges of $ \ccclass $ are incident on $ a $, namely, $ \ccclass $ is a star graph with $ a $ as the center.
    Hence, deleting $ D $ results in a graph with one component containing $ a $ and all other components are singletons, and there are at most $ \card{D} $ of them.

    $ \bullet $ If $ \card{\ccclass \cap \anchorSet} > 1 $, consider two vertices $ a \neq b \in \ccclass \cap \anchorSet $. We show that the number of edge-disjoint paths of length at most $ 2 $ between $ a $ and $ b $ in $ \sparseG[\ccclass] $ is more than $ \card{D} $.

    Note that, for any two vertices $ w, w' \in \openN{\sparseG}{a} \cap \openN{\sparseG}{b} $, the paths $ (a,w,b) $ and $ (a,w',b) $ are of length $ 2 $ and edge-disjoint.
    If $ (a,b) $ is not an edge in $ \sparseG $, the number of such that is $ \card{\openN{\sparseG}{a} \cap \openN{\sparseG}{b}} \geq 2(\frac{3}{4}\card{\ccclass}-1) - (\card{\ccclass} - 2) \geq \frac{1}{2}\card{C} > \card{D} $.
    If $ (a,b) $ is an edge in $ \sparseG $, then $ \card{\openN{\sparseG}{a} \cap \openN{\sparseG}{b}} + \geq 2(\frac{3}{4}\card{\ccclass}-1) - \card{\ccclass} \geq \frac{1}{2}\card{C} - 2 $.
    The number of edge-disjoint paths of length at most $ 2 $ between $ a $ and $ b $, which also includes $ (a,b) $ in this case, is $ \card{\openN{\sparseG}{a} \cap \openN{\sparseG}{b}} + 1 \geq \frac{1}{2}\card{\ccclass} - 1 > \card{D} $.
    Hence, after deleting the edges in $ D $, there is still a path between $ a $ and $ b $.
    As this holds for any two vertices in $ \ccclass \cap \anchorSet $, all vertices in $ \ccclass \cap \anchorSet $ are in the same connected component of $ \sparseG[\ccclass]\setminus D $.

    After deleting $ D $, for any vertex in $ \ccclass \setminus \anchorSet $ that is not in the same component as $ \ccclass \cap \anchorSet $, all of its incident edges must be in $ D $, as the other endpoint of any such edge is in $ \ccclass \cap \anchorSet $.
    This also implies that all other components are singletons, and there are at most $ \card{D} $ of them.

\end{proof}
\begin{lemma}
\label{lem:labelsync.classes.are.connected.components}
Conditioned on $ \successDA $, after every update $ t $, the classes maintained by $ \LabelSync $ are exactly the connected components of $ \sparseG{t} $.
\end{lemma}
\begin{proof}
    We prove by induction on time step $ t $.

    At the base case $ t=1 $, the lemma holds as both $ \sparseG{1} $ and the classes maintained by $ \LabelSync{1} $ are both $ \{v_{1}\} $.

    For the inductive step, assume that the lemma holds for time step $ t-1 $, that is, every class maintained by $ \LabelSync{t-1} $ is a connected component of $ \sparseG{t-1} $.

    We show that it also holds for time step $ t $.

    $ \bullet $ First, we demonstrate that, after the split step (\cref{line:labelsync.split.begin}-\cref{line:labelsync.split.end}), every vertex in same \piece{} is in the same connected component of $ \sparseG{t} $. 

    Indeed, for each class $ \ccclass $ not participating in the split step ($ d_{\ccclass} = 0 $), they have no deleted edges.
    By induction hypothesis, $ \ccclass $ is a connected component of $ \sparseG{t-1} $. With none of its edges deleted, $ \ccclass $ remains a connected component of $ \sparseG{t} $.
    Hence, all its vertices are in the same connected component of $ \sparseG{t} $.

    It remains to consider classes $ \ccclass $ participating in the split step ($ d_{\ccclass} \geq 1 $).
    If $ d_{\ccclass} \geq \frac{1}{2}\card{\ccclass}-1 $, the algorithm iterates over all the vertices in $ \ccclass $, grouping vertices in the same connected component of $ \sparseG{t} $ together (\cref{line:labelsync.split.large.deletion.begin}-\cref{line:labelsync.split.large.deletion.end}).
    Note that, determining whether two vertices $ x,y $ are in the same connected component of $ \sparseG{t} $ can be done by checking if $ \comp{x} = \comp{y} $. 
    If $ d_{\ccclass} < \frac{1}{2}\card{\ccclass} - 1 $, conditioned on $ \successDA $, by~\cref{lem:ccclass.is.dense}, we have $ \openN{\sparseG{t-1}}{a} \geq \frac{3}{4}\card{\ccclass}-1 $ for any $ a \in \ccclass \cap \anchorSet{t-1} $.
    Now, by \cref{lem:labelsync.robust.components}, all vertices in $ \ccclass \cap \anchorSet{t-1} $ are in the same connected component of $ \sparseG{t} $, and all other vertices in $ \ccclass $ that are not in the same component as $ \ccclass \cap \anchorSet{t-1} $ are singletons.

    $ \bullet $ To finish the proof, it suffices to show that, the merge step (\cref{line:labelsync.merge.begin}-\cref{line:labelsync.merge.end}) merges all \piece{s} that are in the same connected component of $ \sparseG{t} $ and assigns them the same label, and assigns different labels to \piece{s} in different components. 
    This step begins with a set $ \piecesToMerge $, consisting of all the \piece{s} created in the split step, and the singleton $ \{v_t\} $; then adding to $ \piecesToMerge $ all the \piece{s} containing endpoints of an edge in $ \addt{t} $ (\cref{line:labelsync.merge.begin}-\cref{line:labelsync.piecesToMerge.include.addt}).
    All the \piece{s} participating in the merge step are the \piece{s} in $ \piecesToMerge $.

    If a \piece{} does not participate in the merge step, that is, is not included in the set $ \piecesToMerge $, then all of its vertices are not incident to any edge in $ \addt{t} $.
    By induction hypothesis, this \piece{} is a connected component of $ \sparseG{t-1} $.
    As no edges are added to this \piece{}, it remains a connected component of $ \sparseG{t} $ and keeps its label.

    For each \piece{} $ \pieceToMerge \in \piecesToMerge $, \cref{line:labelsync.assign.key} assigns a key $ \mathrm{key}(\pieceToMerge) = \comp{x} $, for some $ x \in \pieceToMerge $.
    We have shown above that all vertices in $ \pieceToMerge $ are in the same connected component of $ \sparseG{t} $, and hence, they have the same component identifier.
    It follows that, two \piece{s} share the same key if and only if they are in the same connected component of $ \sparseG{t} $.
    By grouping \piece{s} by their keys in \cref{line:labelsync.group.key}, the algorithm merges all \piece{s} which forms exactly one connected component of $ \sparseG{t} $.
    These \piece{s} are assigned the same label in \cref{line:labelsync.assign.label}.

    Finally, we argue that different components receive different labels. Indeed, each old label is inherited by exactly one retained \piece{}, and fresh labels are distinct and have never been used before.
    The classes after update $ t $ are therefore exactly the components of $ \sparseG{t} $.
\end{proof}

\begin{lemma}\label{lem:number.of.updated.vertices}
    For every $ t\leq n $, the total number of label changes of arrived vertices after executing $ \LabelSync $ (\cref{alg.label.sync}) for the first $ t $ updates is at most $ t\log_{2}n + 4(\log_{2}n + 3)\sum_{t'\leq t} \card{\delt{t'}} $.
\end{lemma}
\newcommand{\newpiecesize}[1]{\sigma_{#1}}
\newcommand{\smallpiecesize}[1]{\pi_{#1}}
\newcommand{\potential}[1]{\mathrm{Budget}(#1)}
\begin{proof}
    Fix any update time $ t'\in [1,t] $.
    We count the number of label changes after executing $ \LabelSync $ for update $ t' $, that is, the number of vertices with different labels between time $ t'-1 $ and $ t' $.

    At a high level, we upper bound this quantity by the number of vertices participating in the split and merge steps of $ \LabelSync $ for update $ t' $.

    Formally, consider the execution of $ \LabelSync $ for update $ t' $.
    Let $ \newpiecesize{t'} $ be the total size of the new \piece{s} created in the split step (\cref{line:labelsync.split.begin}-\cref{line:labelsync.split.end}). 
    In the merge step (\cref{line:labelsync.merge.begin}-\cref{line:labelsync.merge.end}), denote by $ \smallpiecesize{t'} $ be the total size of all the \piece{s} in $ \piecesToMerge $ that are not $ \pieceToMerge^* $.

    We first claim that, the number of vertices with changed labels is at most $ \newpiecesize{t'} + \smallpiecesize{t'} $. Indeed, the vertices with changed labels must belong to:
    \begin{itemize}
        \item Either one of the \piece{s} created in the split step, which also includes \piece{s} $ \piecesToMerge^{*} $ that requires a fresh label in the merge steps. The total size of these \piece{s} is at most $ \newpiecesize{t'} $.
        \item Or one of the \piece{s} in $ \piecesToMerge $ that are not $ \pieceToMerge^* $, which has total size at most $ \smallpiecesize{t'} $.
    \end{itemize}

    It suffices to demonstrate that $ \sum_{t'\leq t} \newpiecesize{t'} + \sum_{t'\leq t}\smallpiecesize{t'} \leq t \log_{2}n + 4(\log_{2} n + 3) \sum_{t'\leq t}\card{\delt{t'}} $.

    \paragraph{Bounding $ \newpiecesize{t'} $.}\label{par:bounding.newpiecesize} % (fold)
    In the split step, if $ d_{\ccclass} \geq \card{\ccclass}/2-1 $, the algorithm traverses all vertices in $ \ccclass $ and groups them into \piece{s} according to their identifiers. 
    The total size of the new \piece{s} created in this case is therefore at most $ \card{\ccclass} \leq 2d_{\ccclass}+2 \leq 4d_{\ccclass} $.
    In the remaining case, if $ d_{\ccclass} < \card{\ccclass}/2-1 $, the new \piece{s} are created grouping vertices in $ W_{\ccclass} $ according to their identifiers.
    The total size of the new \piece{s} created in this case is at most $ \card{W_{\ccclass}} $.
    By definition of $ W_{\ccclass} $ and the invariant~\eqref{eq:invariant.anchor}, we have $ \card{W_{\ccclass}} \leq d_{\ccclass} $.
    In either cases, the total size of the new \piece{s} created for class $ \ccclass $ is at most $ 4d_{\ccclass} $.
    Thus, the total size of the new \piece{s} created in the split step is
    \begin{align}
        \newpiecesize{t'} \leq \sum_{\ccclass \text{ participated in the split step}} 4d_{\ccclass} = 4\card{\delt{t'}}. \label{eq:bound.newpiecesize}
    \end{align}
    % paragraph Bounding newpiecesize (end)

    \paragraph{Bounding $ \smallpiecesize{t'}.$}\label{par:bounding_smallpiecesize_t_} % (fold)
    We show an upper bound for $ \sum_{t'\leq t} \smallpiecesize{t'} $ instead of bounding $ \smallpiecesize{t'} $ individually for each $ t' $.

    Consider the following scenario: throughout the whole course of the algorithm, for all time step $ t' $, there are no classes participating in the split step.
    Observe the merge step, a vertex $ v $ changes its label only if the \piece{} containing $ v $ joins a \piece{} with a larger size.
    This operation at least doubles the size of the \piece{} containing $ v $.
    Since the split opertation is never performed, it follows that any vertex can change its label at most $ \log_{2}n $ times.

    To account for the split step, we employ a potential function argument, which enables us to track the number of label changes allowed for each vertex
    Intuitively, every newly arrived vertex $ v $ is assigned a potential $ \potential{\{v\}} $ of $ \log_{2}n $, which corresponds to the maximum number of label changes it can undergo if the class containing $ v $ never participates in the split step.
    A split step performed on a class reduces the potential of each vertex in that class, while a merge step increases the potential of each vertex in the classes that get merged into a larger class.

    Formally, we define the potential function as follows. For a partition $ \ccclasses $ of a set $ V $, the potential of a class $ \ccclass\in\ccclasses $ is defined as
    \begin{align}
        \potential{\ccclass} = \card{\ccclass}\log_{2} \frac{n}{\card{\ccclass}}, \nonumber
    \end{align}
    and the potential of the whole partition is
    \begin{align}
        \potential{\ccclasses} = \sum_{\ccclass\in\ccclasses} \potential{\ccclass} = \sum_{\ccclass\in\ccclasses} \card{\ccclass}\log_{2} \frac{n}{\card{\ccclass}}. \nonumber
    \end{align}
    For each newly arrived vertex $ v $, assign a potential $ \potential{\{v\}} = \log_{2}n $.

    We analyze the change in potential of the partition maintained by $ \LabelSync $ at following three points in time: $ t'-1 $, after the split step of update $ t' $, and after the merge step of update $ t' $.
    At time $ t'-1 $, the partition is $ \potential{\ccclasses{t'-1}} $, which is the partition of classes maintained by $ \LabelSync $ at time $ t'-1 $, serving as the input to $ \LabelSync{t'} $. The partition after the merge step of update $ t' $ is the output of $ \LabelSync{t'} $. The partition after the split step of update $ t' $, denoted by $ \widetilde{\ccclasses{t'}} $, consists of $ v_{t'} $, the new \piece{s} created in the split step, and the \piece{s} that are not split.

    The split step transforming $ \ccclasses{t'-1} $ into $ \widetilde{\ccclasses{t'}} $ increases the potential. In particular, for each class $ \ccclass \in \ccclass{t'-1} $ split into $ \pieceToMerge_{1}, \pieceToMerge_{2}, \dots  $, with $ \pieceToMerge_{1} $ being the retained \piece{}, the potential increases by
    \begin{align}
        \sum_{i} \card{\pieceToMerge_{i}}\log_{2} \frac{n}{\card{\pieceToMerge_{i}}} - \card{\ccclass}\log_{2} \frac{n}{\card{\ccclass}} &= \sum_{i}\card{\pieceToMerge_{i}} \log_{2}\dfrac{\card{\ccclass}}{\card{\pieceToMerge_{i}}} \nonumber \\
                                                                                                                                                 &\leq \card{\pieceToMerge_{1}}\log_{2}\left(1+\dfrac{\card{\ccclass}-\card{\pieceToMerge_{1}}}{\card{\pieceToMerge_{1}}}\right) + (\card{\ccclass} - \card{\pieceToMerge_{1}})\log_{2}n  \nonumber \\
                                                                                                                                                 &\leq (\card{\ccclass} - \card{\pieceToMerge_{1}})(\log_{2}e + \log_{2}n). \nonumber
    \end{align}
    Note that, newly arrived vertex $ v_{t'} $ contributes $ \log_{2}n $ to the potential, and $ \sum_{\ccclass\in \ccclass{t'-1}}\card{\ccclass} - \card{\pieceToMerge_{1}} \leq \newpiecesize{t'} $.
    Hence,
    \begin{align}
        \potential{\widetilde{\ccclasses{t'}}} - \potential{\ccclasses{t'-1}} \leq \newpiecesize{t'}(\log_{2}e + \log_{2}n) + \log_{2}n. \label{eq.split.increases.potential}
    \end{align}

    The merge step transforming $ \widetilde{\ccclasses{t'}} $ into $ \ccclasses{t'} $ decreases the potential, and we argue that the decrease is at least $ \smallpiecesize{t'} $. 
    Indeed, consider a merge of $ \pieceToMerge_{1}, \pieceToMerge_{2}, \dots, $ with $ \pieceToMerge^* $ into $ \ccclass' $. Here, $ \card{\pieceToMerge^*} \geq \pieceToMerge_{i} $, and hence $ \card{\ccclass}\geq 2 \card{\pieceToMerge_{i}} $.
    Thus, for each such merge, the potential decreases by at least
    \begin{align}
        \sum_{i} \card{\pieceToMerge_{i}}\log_{2} \frac{n}{\card{\pieceToMerge_{i}}} - \card{\ccclass'}\log_{2} \frac{n}{\card{\ccclass'}} &\geq \sum_{i}\card{\pieceToMerge_{i}} \log_{2}\dfrac{\card{\ccclass'}}{\card{\pieceToMerge_{i}}} \geq \sum_{i} \card{\pieceToMerge_{i}}. \nonumber
    \end{align}
    By definition, $ \smallpiecesize{t'} $ covers the total size of all such $ \pieceToMerge_{i} $.
    It follows that $ \potential{\widetilde{\ccclasses{t'}}} - \potential{\ccclasses{t'}} \geq \smallpiecesize{t'} $.

    Taking the sum over all $ t'\leq t $, and note that $ \potential{\ccclasses{0}} = 0 \leq \potential{\ccclasses{t}} $, we have
    { \small
    \begin{align}
        \sum_{t'\leq t} \smallpiecesize{t'} &\leq (\potential{\widetilde{\ccclasses{1}}} - \potential{\ccclasses{1}}) + (\potential{\widetilde{\ccclasses{2}}} - \potential{\ccclasses{2}}) + \dots + (\potential{\widetilde{\ccclasses{t}}} - \potential{\ccclasses{t}}) \nonumber \\
                                            &\leq (\potential{\widetilde{\ccclasses{1}}} - \potential{\ccclasses{0}}) + (\potential{\widetilde{\ccclasses{2}}} - \potential{\ccclasses{1}}) + \dots + (\potential{\widetilde{\ccclasses{t}}} - \potential{\ccclasses{t-1}}) \nonumber \\
                                            &\leq \sum_{t'\leq t} \newpiecesize{t'}(\log_{2}e + \log_{2}n) + t\log n, \label{eq:bound.smallpiecesize}
    \end{align} 
    }
    where the last inequality follows from~\eqref{eq.split.increases.potential}.

    To finish the proof, combining~\eqref{eq:bound.newpiecesize} and~\eqref{eq:bound.smallpiecesize}, we obtain
    \begin{align}
        \sum_{t'\leq t} \newpiecesize{t'} + \sum_{t'\leq t}\smallpiecesize{t'} &\leq \sum_{t'\leq t} 4\card{\delt{t'}} + \sum_{t'\leq t} \newpiecesize{t'}(\log_{2}e + \log_{2}n) + t\log n \nonumber \\
                                                                                  &\leq 4(\log_{2}n + 3)\sum_{t'\leq t}\card{\delt{t'}} + t\log n. \nonumber
    \end{align}
    % paragraph Bounding $ \smallpiecesize{t'} $ (end)
\end{proof}


\subsection{Auxiliary density bound}
\label{sub:auxiliary_lemmas_from_}

\begin{lemma}
\label{lem:ccclass.is.dense}
Conditioned on $ \successDA $, for every time $ t $, every component $ \ccclass $ of $ \sparseG{t} $ with $ \card{\ccclass}>1 $, and every $ a\in \ccclass\cap\anchorSet{t} $,
\[
  |\openN{\sparseG{t}}{a}|\ge\tfrac34\card{\ccclass}-1.
\]
\end{lemma}

% subsection Auxiliary Lemmas From~ (end)

% section Explicit Dynamic Correlation Clustering Algorithm for Insertion-Only Stream (end)
